\subsection{Clustered spectra and attained prediction laws}
\label{sec:v35-spectrum}
The full spectrum, rather than only its least singular value, is a
natural comparison object. Write
\[
 F_L(x)=(e^{\mathrm i kx_j})_{0\le k\le L,\,1\le j\le q}.
\]
Batenkov, Diederichs, Goldman and Yomdin
\cite[Theorems 2.2--2.3]{BDGGY2021} study this Fourier--Vandermonde
matrix on the circle. For an $s$-point cluster with all pairwise
circular distances between $\tau h$ and $h$, their single-cluster
theorem gives
\[
 \sigma_j(F_L)\asymp L^{1/2}(Lh)^{j-1},\qquad 1\le j\le s,
\]
when $L\ge s$ and $Lh$ is sufficiently small, with constants depending
on $s,\tau$. Their multi-cluster theorem compares the complete spectrum
with the ordered union of the cluster spectra, with multiplicative
errors controlled by $(L\vartheta)^{-1}+Lh$, where $\vartheta$ is
inter-cluster separation and $h$ the largest cluster diameter, in the
stated range $C/\vartheta\le L\le c/h$. Thus the spectral comparison
is not confined to the worst-conditioned direction.

For our fixed horizon, the connection to those matrix objects admits
a short direct proof. It includes exact repeated labels and requires
no choice of a collision path.

\begin{proposition}[Exterior spectral products and determinant volumes]
\label{prop:v35-exterior-spectrum}
Fix integers $q\ge1$ and $L\ge q-1$. For $x\in[0,1]^q$, with
repeated entries permitted, let $V_\ell(x)$ be the largest absolute
Vandermonde product on $\ell$ entries, with $V_1=1$. If
$\sigma_1\ge\cdots\ge\sigma_q\ge0$ are the singular values of
$F_L(x)$, there are constants $c_q,C_{q,L}>0$ such that
\begin{equation}\label{eq:v35-spectral-volume}
 c_q V_\ell(x)\le\prod_{j=1}^{\ell}\sigma_j
          \le C_{q,L}V_\ell(x),\qquad 1\le\ell\le q.
\end{equation}
In particular both sides vanish exactly when fewer than $\ell$
distinct node values remain. The constants are independent of all
node gaps; $q$ and $L$ are fixed.
\end{proposition}
\begin{proof}
Put $z_j=e^{\mathrm i x_j}$. For $|x_j-x_i|\le1$,
\[
 2\sin(1/2)|x_j-x_i|\le |z_j-z_i|\le |x_j-x_i|.
\]
A minor with columns $j_1,\ldots,j_\ell$ and rows
$0\le k_1<\cdots<k_\ell\le L$ is an alternating polynomial
$\det(z_{j_b}^{k_a})_{a,b=1}^{\ell}$. It is divisible by every
$z_{j_b}-z_{j_a}$, and hence by their product: these distinct linear
factors are relatively prime in the polynomial ring. The quotient
is a polynomial. There are finitely many row choices, so their
quotients have a common bound on the unit torus, depending only on
$\ell,L$. Every $\ell$-minor is therefore bounded by a constant
times $V_\ell(x)$. This polynomial factorization remains valid
at repeated nodes without division by a numerical gap.

The matrix of $\bigwedge^\ell F_L$ in the standard exterior bases
consists of these minors. Its operator norm is bounded by its
Frobenius norm, hence by
$\binom{L+1}{\ell}^{1/2}\binom q\ell^{1/2}$ times their largest
absolute value. Conversely, take columns attaining $V_\ell$ and
the first $\ell$ rows. Their minor is precisely the Vandermonde
product in the corresponding $z$'s, so its absolute value is at
least $(2\sin(1/2))^{\ell(\ell-1)/2}V_\ell(x)$.
The operator norm dominates any matrix entry. Finally, the singular
value decomposition gives
$\|\bigwedge^\ell F_L\|=\prod_{j\le\ell}\sigma_j$.
Taking the minimum and maximum of the finitely many constants proves
\eqref{eq:v35-spectral-volume}, including all zero-rank cases.
\end{proof}

This proposition uses elementary exterior algebra and polynomial
factorization; it is a comparison identity, not a new clustered-matrix
estimate with uniform growing-bandwidth constants. In particular it
does not replace the dependence on $L$ established in the cited work.

\begin{corollary}[Spectral expression of the acquired envelope]
\label{cor:v35-spectral-envelope}
For each $1\le m<N$, fix an integer $L_m\ge q_m-1$ and apply
$F_{L_m}$ to the formal positive nodes $x_\alpha(a)$ of
Section~\ref{sec:collision-geometry}. Write its singular values as
$\sigma_{m,1}(a)\ge\cdots\ge\sigma_{m,q_m}(a)$. Then
\begin{equation}\label{eq:v35-spectral-envelope}
 \Psi_{n,m}(M,a)\asymp
 \max_{1\le\ell\le p_{n,m}}
    \left(\frac{\prod_{j=1}^{\ell}\sigma_{m,j}(a)}{M}\right)^{2/\ell}.
\end{equation}
The comparison is uniform in every integer $M\ge1$ and $a\in K$,
including exact collisions. The same spectral envelope, maximized
over $n+m=N$, is comparable to either common-controller risk in
Theorem~\ref{thm:resolution-main}.
\end{corollary}
\begin{proof}
Use Proposition~\ref{prop:v35-exterior-spectrum} with
$q=q_m$ and $V_\ell=\mathcal V_{m,\ell}$. Raise its inequalities
to $2/\ell$ and maximize over the finite set of allowed indices.
The last assertion is Theorem~\ref{thm:resolution-main}.
\end{proof}

The matrix in this corollary is an auxiliary function of the known
exponent labels, not an extra observation supplied to the predictor.
Its full spectrum describes the available collision scales. The
upper endpoint $p_{n,m}=\min\{n(r-1),q_m\}$ and the interpretation
of those scales as prediction regret require the acquired tangent,
probability minorization, whole-image cover and causal update proved
in the preceding sections. The following example separates these
roles while holding the future spectrum fixed.

\begin{corollary}[Dependence on acquisition at a fixed future test space]
\label{cor:v35-fixed-future}
Use the detector, prior and command cube of
Corollary~\ref{cor:two-parameter}, and keep the same two-trial query
menu. For any fixed $L\ge8$, the nine singular values of the
auxiliary matrix on its positive future nodes satisfy, componentwise,
\begin{equation}\label{eq:v35-two-parameter-spectrum}
 (\sigma_{2,1},\ldots,\sigma_{2,9})
       \asymp(1,1,1,1,1,1,\rho,\rho,\tau)
\end{equation}
uniformly on the calibration square. The checkpoint problems
$(n,2)$, $n=1,2,3$, with their respective total horizons $n+2$,
have unconditional optimal regrets
\begin{equation}\label{eq:v35-fixed-future}
 \begin{split}
 \overline R_{M;1,2}^{u,v}&\asymp M^{-2/3},\\
 \overline R_{M;2,2}^{u,v}&\asymp M^{-1/3},\\
 \overline R_{M;3,2}^{u,v}&\asymp
  \max\{M^{-1/3},\rho^{1/2}M^{-1/4},
                       (\rho^2\tau)^{2/9}M^{-2/9}\}.
 \end{split}
\end{equation}
The same comparisons hold for their worst-history checkpoint optima.
\end{corollary}
\begin{proof}
The six separated groups in the proof of
Corollary~\ref{cor:two-parameter} give
$\mathcal V_{2,\ell}\asymp1$ for $\ell\le6$ and successive
volumes of orders $\rho,\rho^2,\rho^2\tau$ for $\ell=7,8,9$.
Proposition~\ref{prop:v35-exterior-spectrum} gives the corresponding
partial spectral products. Dividing successive positive products
gives \eqref{eq:v35-two-parameter-spectrum}. At $\tau=0$ with
$\rho>0$ the last singular value is zero, and at $\rho=0$ the last
three vanish, by the exact-rank clause of the proposition.
The truncations are $p_{1,2}=3$, $p_{2,2}=6$, $p_{3,2}=9$.
Substituting them in Theorem~\ref{thm:intrinsic-checkpoint} gives
the first two laws; the third is the checkpoint computation in
Corollary~\ref{cor:two-parameter}.
\end{proof}

Thus one and the same future spectrum is compatible with different
acquired dimensions and resolution envelopes. This is a consequence
of the classification, not a claim that singular values contain no
useful collision information. Nor do the three checkpoint problems
assert a new sharp price for a multi-checkpoint controller; their
horizons are stated separately.

\subsection{Superresolution minimax loss and retained-label loss}
Batenkov, Demanet, Goldman and Yomdin
\cite[Definition 3.13 and Theorem 3.14]{BDGY2020} obtain statistical
minimax consequences, not only conditioning estimates. They reconstruct
complex coefficients of $s$ point sources on a grid of spacing
$\Delta$ from the entire noisy Fourier function on $[-\Omega,\Omega]$.
The support is also unknown. For normalized $L^2$ noise at most
$\delta$, the coefficient-$\ell^2$
minimax error has upper and lower estimates on the scale
$\delta\,\mathrm{SRF}^{2b-1}$, with
$\mathrm{SRF}=\pi/(\Delta\Omega)$ and clustering parameter $2\le b\le s$.
The upper statement selects arbitrarily large fixed values of
$\mathrm{SRF}$ and then sufficiently small $\Delta$; the lower
statement uses specified cluster parameters and all sufficiently
large $\mathrm{SRF}$ in its stated grid range. These are the separate
quantifiers of that theorem, not a bound asserted for every bandwidth
and configuration.

Here acquisition is a fixed stochastic experiment, and the latent
parameter is not an unknown coefficient vector on a reconstruction
grid. The loss compares a retained prediction with the full-history
prediction in the same experiment. The integer $M$ constrains the
number of labels between reports; it is neither a Fourier sample
count nor the reciprocal of external observation noise. Even the
worst-history criterion concerns prediction after compression, not
recovery of a sparse signal from perturbed spectral data. Consequently
one cannot replace $\Delta\Omega$ or $\delta$ by $M^{-1}$ in the
superresolution formula to obtain our theorem. The relevant common
mathematics is degeneration of exponential test systems; the conversion
to the particular loss and the attained family is different.

\begin{remark}[Determinant laws beyond finite-power paths]
\label{rem:v35-flat-path}
Theorem~\ref{thm:collision-tree} has an additional hypothesis comparing
each nonidentical pair gap to a finite power of its path parameter.
Neither Theorem~\ref{thm:resolution-main} nor
Corollary~\ref{cor:v35-spectral-envelope} uses that hypothesis.
For example set, for sufficiently small $\theta>0$,
\[
 u=e^{-1/\theta^2},\qquad
 v=e^{-1/\theta^2}+e^{-1/\theta^4},
\]
and extend by $u=v=0$ at zero. Then
$\rho\asymp e^{-1/\theta^2}$ and
$\tau=e^{-1/\theta^4}$. Direct substitution in
\eqref{eq:two-parameter-law} gives the three resolution regimes;
the two crossover budgets have orders
\[
 \rho^{-6}\asymp e^{6/\theta^2},\qquad
 \rho^2\tau^{-8}\asymp e^{8/\theta^4-2/\theta^2}.
\]
They follow by balancing adjacent terms. The relevant gaps are not
comparable to any finite powers of $\theta$, so this example uses
the determinant theorem directly rather than its finite-power tree
specialization.
\end{remark}
